%==============================================================================
% Section 6: Secondary readings of the Euclidean circle
%==============================================================================
\section{Secondary Readings of the Euclidean Circle}
\label{sec:circle}

This section organises two readings of the Euclidean circle $\Sone$.  These
provide a reduced description and a thermal description on top of the
already constructed Euclidean response dictionary.  Crucially, the readings
discussed here do not replace the main dictionary; they supply secondary
interpretations of the same Euclidean data.

\subsection{KK Reduction Reading}
\label{sec:circle:kk}

In the KK reading we distinguish the geometric response itself from the
reduced descriptor that describes it in lower-dimensional language.  The
effective coefficients and parity-odd channels appearing in the
lower-dimensional sector are images of the dictionary already constructed,
not the primary objects themselves.

We also fix the geometric setup used in the KK reduction normalisation
identity.  When extracting coefficients we treat the Euclidean circle as a
product/collar direction of length $\beta$ and apply the CZ identity
$\mathrm{NY} = \dd\Ctor$ together with Stokes' theorem on the relative
cobordism
\begin{equation}
  X^4 = M^3 \times I_\beta, \qquad I_\beta = [0, \beta].
\end{equation}
The 3D effective coefficient is defined as a functional on a single physical
boundary slice $Y = M^3$ that absorbs the volume factor along $I_\beta$, so
that here $\int_{\partial X^4}\Ctor$ refers to this single-copy boundary
contribution.  The coefficient $k_{3D}^{\mathrm{tor\text{-}CS}}$ is therefore
the 3D boundary coefficient obtained under inflow/collar normalisation, not
a two-endpoint difference along a closed thermal trace.  Where the
cancellation of the non-zero KK tower is invoked, we assume a real bosonic
KK decomposition on an untwisted product circle.

For this reduced-side reading the following three consistency statements are
useful:
\begin{equation}
  k_{3D}^{\mathrm{tor\text{-}CS}} = \frac{1}{2}\,k_{4D}^{\NY},
  \qquad
  Q_{\mathrm{best}} := \frac{\eta(0)}{\Delta k_{\mathrm{matter}}} = 1,
  \qquad
  k_q \in \mathbb{Z}.
\end{equation}
The first two give normalisation-free statements about the
torsional / matter-coupled parity-odd descriptors, and the last gives the
integrality condition for the reduced Chern--Simons level.

The first identity $k_{3D}^{\mathrm{tor\text{-}CS}} = (1/2)k_{4D}^{\NY}$ is
obtained algebraically as the 3D torsion-CS level after KK reduction along
$\Sone$ of the 4D Nieh--Yan action
$S_\NY = (k_{4D}^{\NY}/8\pi^2)\int \dd\Ctor$
(KK reduction normalisation identity; Appendix~\ref{app:E:E11};
verification: \texttt{proofs/kk\_normalization.py}).  Here, on an untwisted
product circle, the Fourier modes $n$ and $-n$ form complex-conjugate pairs
and the even weights are determined by $n^2$ alone, so that the parity-odd
residue takes the form $n\,w(n^2)$.  The non-zero tower therefore cancels
pairwise and leaves no further parity-odd shift.  The second identity
$Q_{\mathrm{best}} = 1$ follows from the comparison between the spinor-CS
and the APS spectral entry: in $\Nilthree$ one has $\eta(0) = 1$ and
$\Delta k_{\mathrm{matter}} = 1$ (matter-minimal inflow audit;
Appendix~\ref{app:E:E12}).  The third statement $k_q\in\mathbb{Z}$ is the
formal quantisation condition required by large gauge transformations
on the frame bundle, $\pi_3(SO(3)) = \mathbb{Z}$ (formal quantisation on the
frame bundle; Appendix~\ref{app:E:E13}).  These are not new observable
families of the second layer but consistency statements that arise when the
Euclidean response is read in lower-dimensional language.

On the bulk side, the universal Maxwell coefficient
\begin{equation}
  \KMxw = -\frac{L^2}{2r_0^4}
\end{equation}
is reread as a topology-independent kinetic baseline, while the activated
Chern--Simons directions, the uniaxial Higgsing of $\Nilthree$, and the
triplet degeneracy of $\Sthree$ are interpreted as topology-dependent
channels of the reduced vector sector.  The reduced language is therefore a
re-reading of the bulk inventory in lower-dimensional form, not something
prior to the bulk inventory.

On the boundary side, the torsional inflow action, the APS spectral entry,
and the frame-bundle-normalised torsional charge can likewise be mapped to
parity-odd descriptors in lower-dimensional language.  However, $\SCZ$ is
written in terms of $\Ctor$ while $\etaAPS$ belongs to Levi--Civita
spectral data, so the observable-family distinction is preserved even after
projecting onto reduced language.  In particular,
$\etaAPS(\Nilthree) = 1/2$ and $\Ntop(\Sthree) = 6 r_0^2$ both admit
reduced parity-odd descriptors but are not the same quantity.  We do not
claim a complete identification between the reduced-side statements and the
native geometric quantities.

In this sense the KK language does not replace the response dictionary; it
provides a reduced reading of it.  The coefficient correspondences and the
normalisation chain are collected in Appendix~\ref{app:D}.

\subsection{Matsubara-Style Thermal Reading}
\label{sec:circle:matsubara}

In order to overlay a thermal-descriptor reading on the same Euclidean
circle, we adopt
\begin{equation}
  \beta = L
\end{equation}
and describe a Matsubara-style thermal reading as the thermal interpretation.
Here $L$ is the same length of the Euclidean circle that already appears in
the geometric formulation; the thermal reading does not introduce a new
circle but assigns thermal meaning to the existing Euclidean data.

Under this interpretation, Maxwell universality, the CS direction-count law,
the uniaxial Higgsing of $\Nilthree$, and the triplet degeneracy of
$\Sthree$ remain Euclidean structural entries; only the thermal-descriptor
language is added to them.  Moreover, the Matsubara reading does not break
the original product structure $M^3\times \Sone$, so it does not by itself
generate a new parity-odd source.  In this sense, what we obtain is a
$P=0$-protected thermal sector, not a full Lorentzian continuation.

The KK reading and the Matsubara-style thermal reading coexist as two
secondary readings of the same Euclidean circle.  Coexistence on the same
circle, however, does not imply theoretical identity: the former is a
lower-dimensional descriptor language and the latter is a thermal
descriptor language, and the reduced gauge coefficient, boundary inflow
action, phase labels, etc.\ merely admit both readings.

The thermal reading of this section is deliberately confined to the
Euclidean scope.  We do not claim a derivation of full finite-temperature
field theory, a strict construction of the partition function, a detailed
analysis of the Matsubara sum, or real-time transport.  The KK--Matsubara
coexistence table is collected in Appendix~\ref{app:C}, and the descriptor
separation and normalisation chain in Appendix~\ref{app:D}.
